\documentclass[12pt,a4paper]{ctexbook}
\usepackage{geometry}
\geometry{margin=2.2cm}
\usepackage{setspace}
\setstretch{1.35}
\usepackage{hyperref}
\title{全唐诗选·poet.tang.2000}
\author{古典文献整理}
\date{}
\begin{document}
\maketitle
\tableofcontents
\newpage
\section{雜曲歌辭 伊州 歌第三}
\textbf{蓋嘉運}\par\vspace{0.5em}
聞道黃花戍，頻年不解兵。\par
可憐閨裏月，偏照漢家營。\par

\section{雜曲歌辭 伊州 歌第四}
\textbf{蓋嘉運}\par\vspace{0.5em}
千里東歸客，無心憶舊遊。\par
挂帆游白水，高枕到青州。\par

\section{雜曲歌辭 伊州 歌第五}
\textbf{蓋嘉運}\par\vspace{0.5em}
桂殿江烏對，彫屏海燕重。\par
秖應多釀酒，醉罷樂高鐘。\par

\section{雜曲歌辭 伊州 入破第一}
\textbf{蓋嘉運}\par\vspace{0.5em}
千門今夜曉初晴，萬里天河徹帝京。\par
璨璨繁星駕秋色，稜稜霜氣韻鐘聲。\par

\section{雜曲歌辭 伊州 入破第二}
\textbf{蓋嘉運}\par\vspace{0.5em}
長安二月柳依依，西出流沙路漸微。\par
閼氏山上春光少，相府庭邊驛使稀。\par

\section{雜曲歌辭 伊州 入破第三}
\textbf{蓋嘉運}\par\vspace{0.5em}
三秋大漠冷溪山，八月嚴霜變草顏。\par
卷斾風行宵渡磧，銜枚電掃曉應還。\par

\section{雜曲歌辭 伊州 入破第四}
\textbf{蓋嘉運}\par\vspace{0.5em}
行樂三陽早，芳菲二月春。\par
閨中紅粉態，陌上看花人。\par

\section{雜曲歌辭 伊州 入破第五}
\textbf{蓋嘉運}\par\vspace{0.5em}
君住孤山下，煙深夜徑長。\par
轅門渡綠水，遊苑遶垂楊。\par

\section{雜曲歌辭 陸州 歌第一}
\textbf{不詳}\par\vspace{0.5em}
分野中峰變，陰晴衆壑殊。\par
欲投人處宿，隔浦問樵夫。\par

\section{雜曲歌辭 陸州 歌第二}
\textbf{不詳}\par\vspace{0.5em}
共得煙霞徑，東歸山水遊。\par
蕭蕭望林夜，寂寂坐中秋。\par

\section{雜曲歌辭 陸州 歌第三}
\textbf{不詳}\par\vspace{0.5em}
香氣傳空滿，妝花映薄紅。\par
歌聲天仗外，舞態御樓中。\par

\section{雜曲歌辭 陸州 排遍第一}
\textbf{不詳}\par\vspace{0.5em}
樹發花如錦，鶯嗁柳若絲。\par
更逢歡宴地，愁見別離時。\par

\section{雜曲歌辭 陸州 排遍第二}
\textbf{不詳}\par\vspace{0.5em}
明月照秋葉，西風響夜砧。\par
彊言徒自亂，往事不堪尋。\par

\section{雜曲歌辭 陸州 排遍第三}
\textbf{不詳}\par\vspace{0.5em}
坐對銀釭曉，停留玉筯痕。\par
君門常不見，無處謝前恩。\par

\section{雜曲歌辭 陸州 排遍第四}
\textbf{不詳}\par\vspace{0.5em}
曙月當窗滿，征人出塞遙。\par
畫樓終日閉，清管爲誰調。\par

\section{雜曲歌辭 簇拍陸州}
\textbf{不詳}\par\vspace{0.5em}
西去輪臺萬里餘，故鄉音耗日應疎。\par
隴山鸚鵡能言語，爲報閨人數寄書。\par

\section{雜曲歌辭 石州}
\textbf{不詳}\par\vspace{0.5em}
自從君去遠巡邊，終日羅幃獨自眠。\par
看花情轉切，攬鏡淚如泉。\par
一自離君後，嗁多雙臉穿。\par
何時狂虜滅，免得更留連。\par

\section{雜曲歌辭 蓋羅縫 一}
\textbf{不詳}\par\vspace{0.5em}
秦時明月漢時關，萬里征人尚未還。\par
但願龍庭神將在，不教胡馬渡陰山。\par

\section{雜曲歌辭 蓋羅縫 二}
\textbf{不詳}\par\vspace{0.5em}
音書杜絕白狼西，桃李無顏黃鳥嗁。\par
寒鳥春深歸去盡，出門腸斷草萋萋。\par

\section{雜曲歌辭 雙帶子}
\textbf{不詳}\par\vspace{0.5em}
私言切語誰人會，海燕雙飛繞畫梁。\par
君學秋胡不相識，妾亦無心去采桑。\par

\section{雜曲歌辭 崑崙子}
\textbf{不詳}\par\vspace{0.5em}
[揚]子譚經去，淮王載酒過。\par
醉來嗁鳥喚，坐久落花多。\par

\section{雜曲歌辭 祓禊曲 一}
\textbf{不詳}\par\vspace{0.5em}
昨見春條綠，那知秋葉黃。\par
蟬聲猶未斷，寒鴈已成行。\par

\section{雜曲歌辭 祓禊曲 二}
\textbf{不詳}\par\vspace{0.5em}
金谷園中柳，春來已舞腰。\par
那堪好風景，獨上洛陽橋。\par

\section{雜曲歌辭 祓禊曲 三}
\textbf{不詳}\par\vspace{0.5em}
何處堪愁思，花間長樂宮。\par
君王不重客，泣淚向春風。\par

\section{雜曲歌辭 上巳樂}
\textbf{張祜}\par\vspace{0.5em}
猩猩血綵繫頭標，天上齊聲舉畫橈。\par
卻是內人爭意切，六宮羅袖一時招。\par

\section{雜曲歌辭 穆護砂}
\textbf{不詳}\par\vspace{0.5em}
玉管朝朝弄，清歌日日新。\par
折花當驛路，寄與隴頭人。\par

\section{雜曲歌辭 思歸樂 一}
\textbf{不詳}\par\vspace{0.5em}
晚日催弦管，春風入綺羅。\par
杏花如有意，偏落舞衫多。\par

\section{雜曲歌辭 思歸樂 二}
\textbf{不詳}\par\vspace{0.5em}
萬里春應盡，三江鴈亦稀。\par
連天漢水廣，孤客未言歸。\par

\section{雜曲歌辭 金殿樂}
\textbf{不詳}\par\vspace{0.5em}
入夜秋砧動，千門起四鄰。\par
不緣樓上月，應爲隴頭人。\par

\section{雜曲歌辭 胡渭州 一}
\textbf{不詳}\par\vspace{0.5em}
亭亭孤月照行舟，寂寂長江萬里流。\par
鄉國不知何處是？雲山漫漫使人愁。\par

\section{雜曲歌辭 胡渭州 二}
\textbf{不詳}\par\vspace{0.5em}
楊柳千尋色，桃花一苑芳。\par
風吹入簾裏，惟有惹衣香。\par

\section{雜曲歌辭 戎渾}
\textbf{不詳}\par\vspace{0.5em}
風勁角弓鳴，將軍獵渭城。\par
草枯鷹眼疾，雪盡馬蹏輕。\par

\section{雜曲歌辭 牆頭花 一}
\textbf{不詳}\par\vspace{0.5em}
蟋蟀鳴洞房，梧桐落金井。\par
爲君裁舞衣，天寒翦刀冷。\par

\section{雜曲歌辭 牆頭花 二}
\textbf{不詳}\par\vspace{0.5em}
妾有羅衣裳，秦王在時作。\par
爲舞春風多，秋來不堪著。\par

\section{雜曲歌辭 采桑}
\textbf{不詳}\par\vspace{0.5em}
自古多征戰，由來尚甲兵。\par
長驅千里去，一舉兩蕃平。\par
按劒從沙漠歌謠滿帝京寄言天下將須立武功名。\par

\section{雜曲歌辭 [楊]下採桑}
\textbf{不詳}\par\vspace{0.5em}
飛絲惹綠塵，軟葉對孤輪。\par
今朝入園去，物色彊看人。\par

\section{雜曲歌辭 破陣樂 一}
\textbf{不詳}\par\vspace{0.5em}
秋來四面足風沙，塞外征人暫別家。\par
千里不辭行路遠，時光早晚到天涯。\par

\section{雜曲歌辭 破陣樂 二}
\textbf{不詳}\par\vspace{0.5em}
漢兵出頓金微，照日明光鐵衣。\par
百里火幡熖熖，千行雲騎騑騑。\par
蹙踏遼河自竭，鼓噪燕山可飛。\par
正屬四方朝賀，端知萬舞皇威。\par

\section{雜曲歌辭 破陣樂 三}
\textbf{不詳}\par\vspace{0.5em}
少年膽氣凌雲，共許驍雄出羣。\par
匹馬城南挑戰，單刀薊北從軍。\par
一鼓鮮卑送欵，五餌單于解紛。\par
誓欲成名報國，羞將開口論勳。\par

\section{雜曲歌辭 戰勝樂}
\textbf{不詳}\par\vspace{0.5em}
百戰得功名，天兵意氣生。\par
三邊永不戰，此是我皇英。\par

\section{雜曲歌辭 劒南臣}
\textbf{不詳}\par\vspace{0.5em}
不分君恩斷，觀妝視鏡中。\par
容華尚春日，嬌愛已秋風。\par
枕席臨窗曉，屏帷對月空。\par
年年後庭樹，芳悴在深宮。\par

\section{雜曲歌辭 征步郎}
\textbf{不詳}\par\vspace{0.5em}
塞外虜塵飛，頻年度磧西。\par
死生隨玉劒，辛苦向金微。\par

\section{雜曲歌辭 [歎]疆場}
\textbf{不詳}\par\vspace{0.5em}
聞道行人至，妝梳對鏡臺。\par
淚痕猶尚在，笑靨自然開。\par

\section{雜曲歌辭 塞姑}
\textbf{不詳}\par\vspace{0.5em}
昨日盧梅塞口，整見諸人鎮守。\par
都護三年不歸，折盡江邊楊柳。\par

\section{雜曲歌辭 水鼓子}
\textbf{不詳}\par\vspace{0.5em}
雕弓白羽獵初回，薄夜牛羊復下來。\par
夢水河邊秋草合，黑山峰外陣雲開。\par

\section{雜曲歌辭 婆羅門}
\textbf{楊敬述進}\par\vspace{0.5em}
迴樂峰前沙似雪，受降城外月如霜。\par
不知何處吹蘆管？一夜征人盡望鄉。\par

\section{雜曲歌辭 浣沙女 一}
\textbf{不詳}\par\vspace{0.5em}
南陌春風早，東鄰去日斜。\par
千花開瑞錦，香撲美人車。\par

\section{雜曲歌辭 浣沙女 二}
\textbf{不詳}\par\vspace{0.5em}
長樂青門外，宜春小苑東。\par
樓開萬戶上，人向百花中。\par

\section{雜曲歌辭 鎮西}
\textbf{不詳}\par\vspace{0.5em}
天邊物色更無春，祗有羊羣與馬羣。\par
誰家營裏吹羌笛，哀怨教人不忍聞。\par

\section{雜曲歌辭 鎮西 二}
\textbf{不詳}\par\vspace{0.5em}
歲去年來拜聖朝，更無山闕對溪橋。\par
九門楊柳渾無半，猶自千條與萬條。\par

\section{雜曲歌辭 回紇}
\textbf{不詳}\par\vspace{0.5em}
曾聞瀚海使難通，幽閨少婦罷裁縫。\par
緬想邊庭征戰苦，誰能對鏡治愁容。\par
久戍人將老，須臾變作白頭翁。\par

\section{雜曲歌辭 長命女}
\textbf{不詳}\par\vspace{0.5em}
雲送關西雨，風傳渭北秋。\par
孤燈然客夢，寒杵擣鄉愁。\par

\section{雜曲歌辭 醉公子}
\textbf{不詳}\par\vspace{0.5em}
昨日春園飲，今朝倒接䍦。\par
誰人扶上馬，不省下樓時。\par

\section{雜曲歌辭 一片子}
\textbf{不詳}\par\vspace{0.5em}
柳色青山映，梨花雪鳥藏。\par
綠窗桃李下，閑坐歎春芳。\par

\section{雜曲歌辭 甘州}
\textbf{不詳}\par\vspace{0.5em}
欲使傳消息，空書意不任。\par
寄君明月鏡，偏照故人心。\par

\section{雜曲歌辭 濮陽女}
\textbf{不詳}\par\vspace{0.5em}
鴈來書不至，月照獨眠房。\par
賤妾多愁思，不堪秋夜長。\par

\section{雜曲歌辭 相府蓮}
\textbf{不詳}\par\vspace{0.5em}
夜聞鄰婦泣，切切有餘哀。\par
即問緣何事，征人戰未迴。\par

\section{雜曲歌辭 蔟拍相府蓮}
\textbf{不詳}\par\vspace{0.5em}
莫以今時寵，寧無舊日恩。\par
看花滿眼淚，不共楚王言。\par
閨燭無人影，羅屏有夢魂。\par
近來音耗絕，終日望應門。\par

\section{雜曲歌辭 離別難}
\textbf{不詳}\par\vspace{0.5em}
此別難重陳，花深復變人。\par
來時梅覆雪，去日柳含春。\par
物候催行客，歸途淑氣新。\par
剡川今已遠，魂夢暗相親。\par

\section{雜曲歌辭 離別難}
\textbf{白居易}\par\vspace{0.5em}
綠楊陌上送行人，馬去車回一望塵。\par
不覺別時紅淚盡，歸來無淚可霑巾。\par

\section{雜曲歌辭 山鷓鴣 一}
\textbf{不詳}\par\vspace{0.5em}
玉關征戍久，空閨人獨愁。\par
寒露溼青苔，別來蓬鬢秋。\par

\section{雜曲歌辭 山鷓鴣 二}
\textbf{不詳}\par\vspace{0.5em}
人坐青樓晚，鶯語百花時。\par
愁人多自老，腸斷君不知。\par

\section{雜曲歌辭 鷓鴣詞}
\textbf{李益}\par\vspace{0.5em}
湘江斑竹枝，錦翼鷓鴣飛。\par
處處湘雲合，郎從何處歸。\par

\section{雜曲歌辭 鷓鴣詞 一}
\textbf{李涉}\par\vspace{0.5em}
湘江煙水深，沙岸隔楓林。\par
何處鷓鴣飛？日斜斑竹陰。\par
二女虛垂淚，三閭枉自沈。\par
惟有鷓鴣鳥，獨傷行客心。\par

\section{雜曲歌辭 鷓鴣詞 二}
\textbf{李涉}\par\vspace{0.5em}
越岡連越井，越鳥更南飛。\par
何處鷓鴣嗁，夕煙東嶺歸。\par
嶺頭行人少，天涯北客稀。\par
鷓鴣嗁別處，相對淚霑衣。\par

\section{雜曲歌辭 樂世}
\textbf{白居易}\par\vspace{0.5em}
管急絲繁拍漸稠，綠腰宛轉曲終頭。\par
誠知樂世聲聲樂，老病人聽未免愁。\par

\section{雜曲歌辭 急[樂世]}
\textbf{白居易}\par\vspace{0.5em}
正抽碧線繡紅羅，忽聽黃鶯斂翠蛾。\par
秋思冬愁春恨望，大都不得意時多。\par

\section{雜曲歌辭 何滿子}
\textbf{白居易}\par\vspace{0.5em}
世傳滿子是人名，臨就刑時曲始成。\par
一曲四詞歌八疊，從頭便是斷腸聲。\par

\section{雜曲歌辭 何滿子}
\textbf{薛逢}\par\vspace{0.5em}
繫馬宮槐老，持杯店菊黃。\par
故交今不見，流恨滿川光。\par

\section{雜曲歌辭 清平調 一}
\textbf{李白}\par\vspace{0.5em}
雲想衣裳花想容，春風拂檻露華濃。\par
若非羣玉山頭見，會向瑤臺月下逢。\par

\section{雜曲歌辭 清平調 二}
\textbf{李白}\par\vspace{0.5em}
一枝紅豔露凝香，雲雨巫山枉斷腸。\par
借問漢宮誰得似，可憐飛燕倚新妝。\par

\section{雜曲歌辭 清平調 三}
\textbf{李白}\par\vspace{0.5em}
名花傾國兩相歡，長得君王帶笑看。\par
解釋春風無限恨，沈香亭北倚闌干。\par

\section{雜曲歌辭 回波樂}
\textbf{李景伯}\par\vspace{0.5em}
回波爾時酒巵，微臣職在箴規。\par
侍宴既過三爵，諠譁竊恐非儀。\par

\section{雜曲歌辭 聖明樂 一}
\textbf{張仲素}\par\vspace{0.5em}
玉帛殊方至，歌鍾比屋聞。\par
華夷今一貫，同賀聖明君。\par

\section{雜曲歌辭 聖明樂 二}
\textbf{張仲素}\par\vspace{0.5em}
九陌祥煙合，千春瑞月明。\par
宮花將苑柳，先發鳳皇城。\par

\section{雜曲歌辭 聖明樂}
\textbf{令狐楚}\par\vspace{0.5em}
海浪恬丹徼，邊塵靜異山。\par
從今萬里外，不復鎖蕭關。\par

\section{雜曲歌辭 大酺樂}
\textbf{不詳}\par\vspace{0.5em}
淚滴珠難盡，容殘玉易銷。\par
儻隨明月去，莫道夢魂遙。\par

\section{雜曲歌辭 大酺樂 一}
\textbf{杜審言}\par\vspace{0.5em}
聖后乘乾日，皇明御曆辰。\par
紫宮初啓坐，蒼璧正臨春。\par
雷雨垂膏澤，金錢賜下人。\par
詔酺歡賞徧，交泰覩惟新。\par

\section{雜曲歌辭 大酺樂 二}
\textbf{杜審言}\par\vspace{0.5em}
毘陵震澤九州通，士女歡娛萬國同。\par
伐鼓撞鐘驚海上，新妝袨服照江東。\par
梅花落處疑殘雪，柳葉開時任好風。\par
大德不官逢道泰，天長地久屬年豐。\par

\section{雜曲歌辭 大酺樂 一}
\textbf{張祜}\par\vspace{0.5em}
車駕東來值太平，大酺三日洛陽城。\par
小兒一伎竿頭絕，天下傳呼萬歲聲。\par

\end{document}